% !TeX program = lualatex-dev
% ===================================================================
% mwe_leg_listas.tex — MWE-2 da camada Leg (crpsp-leg.sty)
%
% Valida o AMBIENTE LegArtigo (dual-mode via \@currenvir) e as
% listas aninhadas de incisos/alíneas/itens:
%   - \begin{LegArtigo}{...}{caput} … \end{LegArtigo} com §§ e
%     incisos>alíneas>itens no corpo → <L>/<LI> aninhados;
%   - a forma PLANA \LegArtigo{...}{...} no mesmo documento
%     (retrocompatibilidade) → parágrafo com \hangindent desde
%     2026-07-11 (era description one-shot por csname — forma que se
%     provou ser o GATILHO do bug de \hsize na quebra de página; ver
%     mwe_list_pagebreak_bug.tex, bissecção de 2026-07-13).
%
% Compilar 2x: lualatex-dev mwe_leg_listas.tex
% ===================================================================

\DocumentMetadata{
    lang        = pt-BR,
    pdfstandard = ua-2,
    pdfversion  = 2.0,
    testphase   = {phase-III,table,firstaid},
    tagging     = on,
}

\documentclass[11pt,a5paper,twoside,openany]{book}

\usepackage{../../book-crpsp_acessivel}

% dc:title no XMP — exigência PDF/UA-2 (veraPDF 8.11.1-t1)
\hypersetup{pdftitle={MWE Leg: listas aninhadas e dual-mode}}

\begin{document}

\mainmatter
\chapter{Listas da camada Leg}

\section{Norma com estrutura aninhada}

\begin{LegNorma}{urn:lex:br:teste:lei:2026;2}{LEI DE TESTE Nº 2, DE 2026}

\LegEmenta{Testa aninhamento artigo > incisos > alíneas > itens.}

\LegAncora{urn:lex:br:teste:lei:2026;2!art1}%
\begin{LegArtigo}{Art. 1º}{Caput do artigo com estrutura aninhada,
listando os requisitos seguintes:}

\begin{LegIncisos}
    \LegInciso{I}{primeiro inciso do caput;}
    \LegInciso{II}{segundo inciso, com alíneas:}
    \begin{LegAlineas}
        \LegAlinea{a}{primeira alínea do inciso II;}
        \LegAlinea{b}{segunda alínea, com itens:}
        \begin{LegItens}
            \LegItem{1}{primeiro item da alínea b;}
            \LegItem{2}{segundo item da alínea b.}
        \end{LegItens}
    \end{LegAlineas}
    \LegInciso{III}{terceiro inciso, após as alíneas.}
\end{LegIncisos}

\LegParagrafo{§ 1º}{Parágrafo dentro do ambiente do artigo, após os
incisos do caput.}
\LegNotaAlteracao{(Redação dada pela Lei de Teste nº 3, de 2026)}{urn:lex:br:teste:lei:2026;3}

\LegParagrafo{§ 2º}{Segundo parágrafo, seguido de incisos próprios
como irmãos (contrato do gerador):}

\begin{LegIncisos}
    \LegInciso{I}{inciso do § 2º;}
    \LegInciso{II}{\LegRevogado{(Revogado) texto rebaixado
        visualmente, nunca riscado.}}
\end{LegIncisos}

\end{LegArtigo}

\LegAncora{urn:lex:br:teste:lei:2026;2!art2}%
\LegArtigo{Art. 2º}{Artigo na forma PLANA retrocompatível, sem
filhos, no mesmo documento da forma ambiente. Remissão interna:
\LegRef{urn:lex:br:teste:lei:2026;2!art1}{art. 1º desta Lei} e
remissão externa: \LegRef{urn:lex:br:federal:lei:2015-07-06;13146}{Lei
nº 13.146, de 2015}.}

\LegDispDecimal{28.1}{Item decimal de primeiro nível.}
\LegDispDecimal{28.1.1}{Item decimal de segundo nível, com recuo.}
\LegDispDecimal{28.1.1.1}{Item decimal de terceiro nível.}

\LegFecho{Brasília, 10 de julho de 2026; 205º da Independência e
138º da República.}

\LegAssina{NOME DE TESTE}{Conselheira-Presidenta}

\end{LegNorma}

\LegTitulo{Título legado (deve avisar deprecação no log, uma vez)}

\end{document}
